\section{惯性、参考系与运动定律}
\subsection{运动状态怎样由相互作用改变}
运动学给出位置、速度和加速度之间的联系，却不回答某个物体为什么具有这样的加速度。动力学把运动的变化同相互作用联系起来。描述一个正在运动的物体，不能凭速度方向添画一个“运动力”；水平光滑面上的物体即使没有水平外力，也能继续运动。力的作用对应速度矢量的变化，而非速度本身。

\fig{n1deep:motion}{速度相同的物体，可以有不同的加速度}{\begin{tikzpicture}
\foreach \x/\lab in {0/甲,3.4/乙,6.8/丙}{\begin{scope}[xshift=\x cm]\fill(0,0)circle(2pt);\draw[flow,blue](0,.5)--(1.3,.5)node[right]{$\bm v$};\node at(.5,-1){\lab};\end{scope}}
\node at(.6,-.4){$\sum\bm F=0$};\draw[force,red](3.4,0)--(4.7,0)node[right]{$\bm F$};\draw[force,red](6.8,0)--(6.8,-.7)node[left]{$\bm F$};
\end{tikzpicture}}
图\ref{n1deep:motion}中的三个速度箭头完全相同，受力却不同。甲在随后短时间仍保持原速度；乙的速度沿原方向增大；丙的速度获得向下的增量，运动方向开始改变。三者都满足 $\Delta\bm v\approx(\sum\bm F/m)\Delta t$，其中力决定的是增量箭头。由此可以理解：仅知道“向右运动”，既不能判断合力向右，也不能判断合力非零。

牛顿第一定律指出，存在这样的参考系：不受外力的质点保持静止或匀速直线运动。这样的参考系称为惯性参考系。第一定律既确定了惯性运动，也给出第二定律适用的参考系条件；不能仅把它当作在任意坐标系内令合力为零的代数结果。相对于某惯性系作匀速直线平移的参考系仍是惯性系，因为两系测得的加速度相同。

地面在常见短时、局部实验中可近似看作惯性系。这是精度与尺度的近似，不代表地球没有自转。若选取正在加速的车厢，静止在车内的小球可以受到不为零的真实合力；此时不能把车内相对静止直接译成真实合力为零。第二讲将通过参考系变换建立加速系中的方程，本讲先在地面惯性系内计算。

\fig{n1t-frame}{同一小球在地面和车厢中的运动描述}{\begin{tikzpicture}
\draw[world](-2.8,0)--(3,0);\draw[world](-1.8,.3)rectangle(1.8,2.6);\draw[world](.5,2.6)--(0,1.2);\fill(0,1.2)circle(3pt)node[below left]{$m$};
\draw[force,blue](0,1.2)--(.5,2.6)node[right]{$T$};\draw[force,red](0,1.2)--(0,.35)node[right]{$mg$};
\draw[flow](1.9,1.5)--(3,1.5)node[right]{$\bm a_0$};\node at(-1,2){车内静止};\node at(0,-.45){地面系：$\bm a=\bm a_0$};
\end{tikzpicture}}
图\ref{n1t-frame}中，小球与车厢的位置差不变，所以它们相对地面的加速度相同。张力的竖直分量抵消重力，水平分量提供小球的水平加速度。研究对象是小球时，车轮与地面的力不直接画进小球的受力图；这些力影响车厢运动，再经悬线改变小球的运动。

\subsection{第二定律的瞬时性与矢量性}
对质量不变的质点，在惯性系中有
\begin{equation}\label{n1eq-newton}
\sum\bm F=m\bm a.
\end{equation}
这里的力都是同一时刻作用在同一物体上的真实外力，加速度也取同一时刻。若力随时间或位置变化，方程仍成立，只是不能再直接使用匀加速运动公式。质量是惯性的量度；在同一合力下，质量越大，加速度越小。

式\eqref{n1eq-newton}是矢量关系。沿所选正交坐标轴分别写成 $\sum F_x=ma_x$、$\sum F_y=ma_y$。某方向速度为零，不必有该方向加速度为零；例如圆周运动最低点的竖直速度为零，竖直加速度却向上。反过来，某方向速度恒定，才可以在该方向令合力为零。选择坐标的目的，是让已知约束直接进入方程，而不是默认所有方向都平衡。

对固定的一组质点，更一般的形式是 $\bm F_{\rm ext}=\dd\bm P/\dd t$。定质量质点的动量为 $\bm p=m\bm v$，才进一步得到 $m\bm a$。有限时间内的 $\Delta\bm P/\Delta t$ 给出平均合力，只有力恒定或取极限时才能同瞬时合力直接等同。若研究对象不断接收或排出物质，还须计算流入、流出物质带来的动量，不能只把一个随时间变化的质量代入 $F=ma$。

\section{研究对象、受力图与相互作用}
\subsection{一张受力图只属于一个研究对象}
受力分析从物体周围的接触与场开始。重力由地球作用；接触支持力沿局部公法线；理想轻绳的张力沿绳并指向离开物体的一端；弹簧力沿弹簧轴线，其方向由伸长或压缩决定。画图时先标明研究对象，再逐一寻找施力物体。所谓“沿斜面向下的重力分力”是重力在坐标轴上的分解，不能同时把重力和两个分力都计入合力。

\fig{n1t-fbd}{斜面受力与分量对应不同层面的描述}{\begin{tikzpicture}
\draw[world](-3,0)--(0,0)--(-3,1.8)--cycle;\node[draw,rotate=-30.9638,anchor=south]at(-1.7,1.02){$A$};
\begin{scope}[shift={(2,1.1)}]\fill(0,0)circle(2pt);\draw[force,red](0,0)--(0,-1.2)node[right]{$mg$};\draw[force,blue](0,0)--(.65,1.1)node[above]{$N$};\draw[force](0,0)--(-1,.6)node[left]{$f$};\draw[axis](0,0)--(1.4,-.84)node[right]{$t$};\draw[guide](0,-1.2)--(.53,-.32);\end{scope}
\end{tikzpicture}}
图\ref{n1t-fbd}中沿斜面向下取 $t$ 轴。图中的 $f$ 只是先选定的正方向，算得负值就说明实际摩擦方向相反；$N$ 对普通无黏附接触则必须非负。坐标分量可以带正负，接触力大小的非负条件却是模型成立的限制，两者不能混同。

第三定律对同一对相互作用写成 $\bm F_{A\to B}=-\bm F_{B\to A}$。这两个力作用在不同物体上，不能在一个物体的受力图中相消。一对平衡力作用在同一物体上，它们是否相等取决于该方向加速度是否为零。书放在加速上升的电梯内时，书与桌面的相互作用仍等大反向，而桌面对书的支持力大于书的重力。

\fig{n1deep:pair}{加速电梯内的书：同一物体的两力与两个物体间的力对}{\begin{tikzpicture}
\begin{scope}
\node[draw,minimum width=1.5cm,minimum height=.7cm](B)at(0,0){书 $B$};\draw[force,blue](B.north)--(0,1.8)node[above]{$\bm N_{D\to B}$};\draw[force,red](B.south)--(0,-1.4)node[below]{$m\bm g$};\draw[flow](1.7,0)--(1.7,1.1)node[above]{$\bm a$};\node at(.5,-2){只隔离书：$N-mg=ma$};
\end{scope}\begin{scope}[xshift=6cm]
\node[draw,minimum width=1.5cm](B)at(0,.6){书 $B$};\node[draw,minimum width=2cm](D)at(0,-.6){桌 $D$};
\draw[force,blue](B.north)--++(0,1)node[above]{$\bm N_{D\to B}$};\draw[force,blue](D.south)--++(0,-1)node[below]{$\bm N_{B\to D}$};\node at(0,-2.6){两箭头属于不同物体};\end{scope}
\end{tikzpicture}}
图\ref{n1deep:pair}左图中两力都作用在书上，合力向上，故 $N>mg$；它们不是作用反作用力。右图只突出书与桌之间的同一接触：即使电梯加速，桌对书与书对桌的两力仍等大反向。第三定律不要求系统静止，第二定律则用来判断同一物体上各力怎样合成。右图未列出各物体的其他外力，不能把它当作完整受力图直接求加速度。

判断两力是不是作用反作用力，最可靠的方法是写出施力物体与受力物体。例如“地球对书的引力”和“桌面对书的支持力”都作用于书，不能组成第三定律力对；前者的反作用是书对地球的引力，后者的反作用是书对桌面的压力。名称中都出现“重力”“支持”或数值碰巧相等，不足以判定两力之间的逻辑关系。

测力装置显示的也不是抽象的“所受合力”。秤盘的压力传感器测接触法向力，悬挂秤测拉力，压电元件可能测某个特定方向的分量。物体处于匀速运动时，合力可以为零而秤读数并不为零。先明确被测相互作用，再隔离被测物体写动力学，是从仪器读数反推运动状态的前提。

\subsection{轻绳、轻滑轮与读数}
取轻绳上一小段，其质量趋于零而加速度有限时，合力也须趋于零。对无分布外力、与接触面无摩擦的理想绳段，两端张力大小相等。这个结论来自该绳段的动力学，并非力可以不加条件地“传递”。有质量的竖直绳、绕粗糙柱面的绳、不同轻绳以及有转动惯量的滑轮，两侧张力都可能不同。

弹簧秤显示的是其敏感元件所受的拉力。把一个定滑轮挂在秤下，若两股绳均竖直且张力为 $T$，滑轮受到的合向下拉力为 $2T$，忽略滑轮自重时秤读数为 $2T$。若两绳夹角为 $\gamma$，对称条件下合力大小改为 $2T\cos(\gamma/2)$；“读数等于两倍张力”不是所有滑轮装置的通式。

\fig{fv:n1:pulley}{弹簧秤承受的是两段绳张力的矢量和}{\begin{tikzpicture}[scale=.9]
\draw[thick](0,0)circle(.6);\draw[thick](-.6,0)--(-.6,-1.8) (.6,0)--(.6,-1.8);
\draw[->,blue,thick](-.6,-.4)--(-.6,-1.5)node[left]{$T$};\draw[->,blue,thick](.6,-.4)--(.6,-1.5)node[right]{$T$};
\draw[->,red,thick](0,.3)--(0,2)node[right]{$F_s$};\node at(3.6,.4){$F_s-2T=0$};\node at(3.6,-.5){轻滑轮的平动方程};
\end{tikzpicture}}
图\ref{fv:n1:pulley}只选滑轮为研究对象，两个向下的 $T$来自绳对轮的作用，向上的 $F_s$来自支架或秤。理想轻滑轮的质量趋零而受力有限，故合力必须趋零，秤读数为 $2T$。若两绳夹角改变，读数应改为两张力的合力大小；绳张力相等并不意味着支承读数永远是它们的代数和。另选绳端物体列方程时，不能再把 $F_s$当成该物体的受力。


\section{整体方程与约束条件}
\subsection{整体法为什么能够消去内力}
对系统内每个定质量质点分别列第二定律，将各式相加。每对内部相互作用按第三定律抵消，得到
\begin{equation}\label{n1eq-system}
\bm F_{\rm ext}=\sum_i m_i\bm a_i=M\bm a_C,
\qquad \bm r_C=\frac{\sum_i m_i\bm r_i}{M}.
\end{equation}
内力不出现在整体合力中，但它们仍决定内部相对运动。式中右端一般是各质点加速度的质量加权和，只有所有物体具有相同加速度时，才能写成 $M\bm a$。两个滑块沿不同方向运动，即使由同一根绳连接，也不能把它们当作沿某一个方向共同平动。


\fig{n1t-system}{接触力在两个隔离方程中成对出现，在整体方程中抵消}{\begin{tikzpicture}
\draw[world](-3,0)--(3,0);\node[draw,minimum width=1.4cm,minimum height=.8cm,anchor=south](A)at(-.7,0){$A$};\node[draw,minimum width=1.4cm,minimum height=.8cm,anchor=south](B)at(.7,0){$B$};
\draw[guide,rounded corners](-1.6,-.2)rectangle(1.6,1);\draw[force](-2.8,.4)--(A.west)node[midway,above]{$F$};
\begin{scope}[yshift=-2cm]\fill(-1,0)circle(1.5pt)node[below]{$A$};\draw[force](-2.5,0)--(-1,0)node[midway,above]{$F$};\draw[force,blue](-1,-.3)--(-2,-.3)node[midway,below]{$N$};\fill(1,0)circle(1.5pt)node[below]{$B$};\draw[force,blue](1,0)--(2,0)node[right]{$N$};\end{scope}
\end{tikzpicture}}
图\ref{n1t-system}的虚线边界包含 $A,B$ 时，两物体之间的接触力成对消去；若只隔离 $B$，接触力就是它的外力。因而求共同加速度常先取整体，求内部张力或压力再隔离其中一个物体。每一步都应重新检查边界，不能把上一步已消去的内力当作不存在。


例如图\ref{n1t-system}的两块在光滑水平面保持接触，水平方向分别有
\[
F-N=m_Aa,\qquad N=m_Ba.
\]
相加消去 $N$，得到 $F=(m_A+m_B)a$；再代回得到 $N=m_BF/(m_A+m_B)>0$。这不仅算出内力，也验证接触能维持。若将外力反向拉走左块，普通接触不能拉住右块，应改用分离模型，不能沿用共同加速度。

还要区分“内力的矢量和为零”和“内力做功之和为零”。一对内力满足 $\bm F_{12}=-\bm F_{21}$，但两受力点位移可能不同，瞬时总功率为 $\bm F_{12}\cdot(\bm v_2-\bm v_1)$，并不必然为零。弹簧推动两球相互远离时，两球都可能获得动能；摩擦使两个表面相对滑动时，机械能也可能减少。这不违反整体动量方程，因为动量与功分别累积力和力乘位移。整体法消去内力不能被用来断言系统内部不发生能量转化。

对放在地面的固定箱体和内部运动物体，竖直整体方程尤其有用。若箱体质量 $M$、内部物体质量 $m$、内部物体向下加速度为 $a$，且除重力和地面支持力外无其他竖直外力，则 $N-(M+m)g=-ma$。所得 $N$ 不必等于总重力。这个结果并不依赖内部具体传力路径，却要求内部物体给定的加速度能由其实际受力实现。

\subsection{绳长、接触与共同运动}
约束把不同物体的运动联系起来。不可伸长轻绳始终绷紧时，总长度恒定；把长度写成坐标函数再求导，才能得到速度和加速度关系。两段竖直绳支承一个动滑轮时，经常有$2y_P-y_1-y_2=\text{常量}$。先对时间求导，得$2\dot y_P-\dot y_1-\dot y_2=0$；再求一次导数，得$2\ddot y_P-\ddot y_1-\ddot y_2=0$，即$2a_P=a_1+a_2$。这里各速度、加速度均为同一坐标正方向的分量。

光滑接触只限制法向相对运动，不意味着切向加速度相同。直斜面作水平平移时，物块与斜面的法向加速度相同；沿斜面则可以相对滑动。若题目说明物块始终相对斜面静止，才可进一步令它们完整的加速度矢量相等。把“接触”和“共同运动”区分开，是处理自由斜面问题的关键。

\fig{n1deep:constraint}{接触只锁定法向分量，相对运动仍可沿切向发生}{\begin{tikzpicture}
\draw[world](-2,-1.2)--(2,1.2);\fill(0,0)circle(2pt)node[above left]{$P$};\draw[axis](0,0)--(1.7,1.02)node[right]{$\bm t$};\draw[axis](0,0)--(-.8,1.3333)node[above]{$\bm n$};
\draw[flow,blue](0,0)--(1.5,0)node[right]{$\bm A$};\draw[flow,red](0,0)--(.5,-.6)node[below left]{$\bm a_m$};\draw[flow,green!50!black](1.5,0)--(.5,-.6)node[midway,below right,yshift=-3mm]{$\bm a_m-\bm A$};
\end{tikzpicture}}
图\ref{n1deep:constraint}以不转动的直斜面为例，蓝箭头是斜面的平移加速度 $\bm A$，红箭头是物块对地加速度。两箭头末端的差矢量平行于斜面，所以接触条件是 $(\bm a_m-\bm A)\cdot\bm n=0$，而不是 $\bm a_m=\bm A$。若物块还相对斜面静止，才要同时令切向相对加速度与初始相对速度为零。曲面或转动表面则还含曲率、法向转动等项，不能照搬这一条直面公式。

普通支持面只能推，绳只能拉。因此假设接触或绷紧后，必须检查 $N\ge0$、$T\ge0$。负的计算值通常意味着假设的约束不能维持，需要改用分离或松绳后的运动方程，而不是将负值取绝对值。

约束还应同初值相容。绳已绷紧且不可伸长时，初始位置必须满足长度关系，初始速度也必须满足其一阶导数；不能任意给两个绳端互不相容的速度再继续按同一模型求有限加速度。如果原来松弛的绳重新绷紧，常需用短时冲击处理速度改变，此时仅靠普通有限力的瞬时方程不足以决定碰后速度。题目若明确忽略这一过程，必须说明后续使用的初速度从何种理想化取得。

\section{摩擦、瞬时变化与临界状态}
\subsection{静摩擦从所需值开始计算}
静摩擦阻碍接触面间的相对滑动趋势，大小由共同运动所需的方程决定，满足 $|f|\le f_{s,\max}$。只有处在临界状态时才令 $|f|=f_{s,\max}$。已经滑动时采用动摩擦 $f_k=\mu_kN$，方向同相对速度相反，不能由物体对地面的速度直接决定。最大静摩擦与动摩擦是否相等，须由题设或明确近似给出。

\fig{n1deep:friction}{缓慢增大的拉力先由静摩擦匹配，超过上限才转入滑动}{\begin{tikzpicture}
\draw[axis](0,0)--(5,0)node[right]{$F$};\draw[axis](0,0)--(0,3.3)node[above]{$|f|$};
\draw[blue,world](0,0)--(2.5,2.5);\fill[blue](2.5,2.5)circle(1.6pt);\draw[guide](2.5,0)--(2.5,2.5)--(0,2.5);\draw[red,world](2.5,1.6)--(4.6,1.6);\draw[red,fill=white](2.5,1.6)circle(1.6pt);\node[left]at(0,2.5){$\mu_smg$};\node[left]at(0,1.6){$\mu_kmg$};\node[above left]at(1.7,1.7){静止};\node[above]at(3.8,1.6){向右滑动};\node[below]at(2.5,0){$F_c$};
\end{tikzpicture}}
图\ref{n1deep:friction}画的是水平面上单个物块，水平向右拉力从零缓慢增大，且 $\mu_k<\mu_s$。静止段由 $F-f=0$ 得到 $f=F$，这就是斜线；到 $F_c=\mu_smg$ 才达到上限。开始向右滑动后，摩擦大小改取 $\mu_kmg$，净力为 $F-\mu_kmg$，因而物块加速。若题设近似 $\mu_k=\mu_s$，图中跳降消失；叠块中还需分别检查每一个接触面的上限，不能把这张单块图的阈值原样套用。

叠块问题应先假设共同运动，求共同加速度和保持上层共同运动所需的摩擦，再检查上限。若超出上限，就分别列滑动方程，同时验证相对加速度确实使预设滑动方向形成。原来静止而拉力逐渐增加的装置，还要先检查地面能否用静摩擦平衡外力；不能直接套用已滑动时的地面摩擦。

\subsection{剪断瞬间，什么可以连续}
理想弹簧服从 $F_s=kx$。无冲量时位置和速度在剪断瞬间连续，所以弹簧的长度、形变和弹力不能在这个瞬间跳变。轻绳张力则是维持约束所需的力，绳被剪断后相应张力立即消失；其他仍绷紧绳段的张力也可能重新分配。于是“弹簧力暂不变”与“绳张力暂不变”是两个不同命题，后者一般错误。

\fig{n1t-cut}{剪断后弹簧形变从原值连续演化，绳张力立即消失}{\begin{tikzpicture}
\draw[world](-2.5,2)--(-.5,2);\draw[spring](-1.5,2)--(-1.5,1);\fill(-1.5,1)circle(3pt)node[left]{$A$};\draw[guide](-1.5,1)--(-1.5,0);\fill(-1.5,0)circle(3pt)node[left]{$B$};\node[right]at(-1.5,.5){剪断};
\begin{scope}[xshift=2cm]\draw[axis](-1.2,0)--(2.3,0)node[right]{$t$};\draw[axis](0,0)--(0,2)node[above]{伸长 $x$};\draw[blue,world](-1.1,1.3)--(0,1.3);\draw[blue,world]plot[domain=0:2,samples=40](\x,{1.3-.18*(\x)^2});\node[left]at(0,1.3){$x_0$};\node[below]at(0,0){$0$};\node at(.5,-.7){初静止时切线也连续};\end{scope}
\end{tikzpicture}}
图\ref{n1t-cut}中，剪断之后两物体可以立刻具有不同加速度，但不能凭空出现有限的速度跳跃。若剩余绳段仍约束两个物体，必须联立其运动约束重求张力。在初速度为零的定长摆绳问题中，径向加速度 $v^2/l$ 恰为零，这一几何条件常用来确定剪断后尚存的绳张力；当初速度不为零时就不能再令径向加速度为零。

分离临界通常由 $N=0$ 给出，滑动临界通常由 $|f|=f_{s,\max}$ 给出，刚离地前还须保留接触阶段的运动关系。先确定哪一个不等式到达边界，再切换方程，才能避免把“压力等于重力”“摩擦恒等于 $\mu mg$”等特殊结果误作普遍规律。

\subsection{从力的突然改变到运动的连续衔接}
加速度由当前受力决定，位置和速度还保留此前运动的结果。若一个质量$m$的物体在$t=0$前后水平合力从$F_-$突变为有限的恒定值$F_+$，先在切换后积分加速度，再积分速度：
\[
v(t)-v(0)=\int_0^t\frac{F_+}{m}\,\dd\tau
=\frac{F_+}{m}[\tau]_0^t=\frac{F_+t}{m},
\]
\[
x(t)-x(0)=\int_0^t\left[v(0)+\frac{F_+\tau}{m}\right]\dd\tau
=\left[v(0)\tau+\frac{F_+\tau^2}{2m}\right]_0^t.
\]
所以速度没有冲量式跳跃，位置为$x(0)+v(0)t+F_+t^2/(2m)$。因此“撤去拉力”并不是“撤去速度”，“剪断连接”也不表示每个物体马上自由落体，仍应逐个检查剩余真实力。

\fig{n1t-step}{有限力阶跃使加速度跳变，速度仍连续}{\begin{tikzpicture}
\begin{scope}\draw[axis](0,0)--(2.6,0)node[right]{$t$};\draw[axis](0,0)--(0,2)node[above]{$F$};\draw[world](0,.5)--(1,.5);\draw[world](1,1.5)--(2.4,1.5);\draw[guide](1,.5)--(1,1.5);\node[below]at(1,0){$0$};\end{scope}
\begin{scope}[shift={(4,0)}]\draw[axis](0,0)--(2.6,0)node[right]{$t$};\draw[axis](0,0)--(0,2)node[above]{$v$};\draw[world](0,.5)--(1,.8)--(2.3,1.8);\draw[guide](1,0)--(1,1.8);\node[below]at(1,0){$0$};\end{scope}
\end{tikzpicture}}
图\ref{n1t-step}把相同初始速度与不同后续斜率区分开。解答应明确所求是切换前极限、切换后瞬间，还是经过一段时间的结果。弹簧力连续属于剪断后瞬间的陈述；如果要问稍后物体是否再次接触，则必须继续积分并检查间隙，不能把瞬时受力图当作永久的过程图。
